\documentclass[12pt, a4paper]{article}

\usepackage[utf8]{inputenc} %Pour l'encodage en utf-8
\usepackage[T1]{fontenc} %Pour les accents
\usepackage[french]{babel} %Pour les langues

\usepackage{siunitx} %Pour les unités
\usepackage{fancyhdr} %Pour les entêtes et bas de page
\usepackage{graphicx} %Pour les figures
\usepackage{mathtools} %Pour les maths
\usepackage[margin=2.5cm]{geometry} %Pour les marges de la page
\usepackage{hyperref} %Pour les liens et les métadonnées

\hypersetup{
	colorlinks=true,
	linkcolor=black,
	urlcolor=blue,
	pdfauthor={Dalina CHABANE, Khalil FALAH, Maxime BELLAUD, Paul ROUSSEAU},
	pdftitle={ATECE - Préparation pour la séance 2}
} %Pour modifier la couleur des liens
\allowdisplaybreaks %Pour avoir une suite d'équations sur plusieurs pages

%Configuration du style de la page
\pagestyle{fancy}
\setlength{\headheight}{14.5pt}
\lhead{ATCE}
\rhead{Préparation pour la séance 2}

\title{ATCE - Préparation pour la séance 2}
\author{Dalina CHABANE, Khalil FALAH, \\ Maxime BELLAUD, Paul ROUSSEAU}
\date{31/01/2025}

\begin{document}

\maketitle

\section{\texorpdfstring{Broches du microcontrôleur qu'il faut absolument \linebreak connecter pour le faire fonctionner, en plus des \linebreak alimentations et du résonateur à quartz}{Broches du microcontrôleur qu'il faut absolument connecter pour le faire fonctionner, en plus des alimentations et du résonateur à quartz}}

\begin{itemize}
	\item Alimentation : Vcc (broches 14 et 34) en \SI{5}{\volt} et GND (broches 15, 23, 35 et 43).
	\item Résonateur : XTAL1 (input, broche 17) et XTAL2 (output, broche 16).
	\item $\overline{\text{RESET}}$ (broche 13), reliée à Vcc via une résistance de tirage.
\end{itemize}

\section{Composants additionnels, en plus du microcontrôleur, que nous sommes susceptibles d'utiliser}

\begin{itemize}
	\item Oscillateur à Quartz de \SI{16}{\mega\hertz}.
	\item Composants passifs (résistances, condensateurs, etc)
	\item Régulateur de tension
	\item Condensateur à valeur variable commandée par une tension (diode varicap)
	\item De quoi téléverser le programme dans le microcontrôleur (USB ?)
	\item Connecteur d'alimentation \SI{5}{\volt} (USB ?)
\end{itemize}

\section{Gestion de la tension d'alimentation du circuit (valeur moyenne et stabilité)}

\begin{itemize}
	\item Pile ou source de tension continue extérieure (\qtyrange{7}{12}{\volt})
	\item Filtrage avec des condensateurs
	\item Diode Zener en protection
	\item Séparation des plans de masse entre les parties puissance/analogique/numérique
	\item Composants déjà faits ? Par exemple le LM7805 ou le LM1117
\end{itemize}

\section{Génération de la tension de commande qui ajustera le condensateur de tirage de l'oscillateur}

\begin{itemize}
	\item Fast PWM avec un filtre passe-bas\footnote{Dans les sections des timers, \textit{Modes of Operation} $\rightarrow$ \textit{Fast PWM} $\rightarrow$ «~This high frequency makes the fast PWM mode well suited for power regulation, rectification, and DAC applications~»}
	\item Utilisation de transistors si une tension supérieure à \SI{5}{\volt} est nécessaire
\end{itemize}

\section{Périphériques du microcontrôleur qui seront nécessaires aux tâches de gestion du temps}

\begin{itemize}
	\item Timers (modes CTC, Output Compare et Input Capture)
	\item Interruptions
\end{itemize}

\end{document}